\textbf{Desarrollo de las magnitudes intermedias.} El punto neutro se obtiene de $0=a_w(h_n-h_{ac})-K_t(h_t-h_n)+C_{m\alpha,f}$. Agrupar $h_n$ da el denominador $a_w+K_t$; después se calcula el margen y la pendiente estática.\begin{enumerate}

\item Pendiente efectiva de la cola.

\[\fitmath{\mathrm{Kt}=0.9\cdot 0.18\cdot 4.2\cdot \left(1-0.35\right)\approx 0.44226\,}\]

\item Resolver el equilibrio de pendiente nula.

\[\fitmath{\mathrm{hn}=\frac{5.1\cdot 0.25+\mathrm{Kt}\cdot 3.1-0.06}{5.1+\mathrm{Kt}}}\]

\[\fitmath{\mathrm{hn}=\frac{5.1\cdot 0.25+(0.44226)\cdot 3.1-0.06}{5.1+(0.44226)}\approx 0.4665977\,}\]

\item Distancia del CG al punto neutro.

\[\fitmath{\mathrm{SM}=\mathrm{hn}-0.35}\]

\[\fitmath{\mathrm{SM}=(0.4665977417)-0.35\approx 0.1165977\,}\]

\item Comprobar tendencia restauradora; la estabilidad dinámica requiere además amortiguamiento y autovalores.

\[\fitmath{\mathrm{Cm}=-\left(5.1+\mathrm{Kt}\right)\cdot \mathrm{SM}}\]

\[\fitmath{\mathrm{Cm}=-\left(5.1+(0.44226)\right)\cdot (0.1165977417)\approx -0.646215\,\mathrm{1/rad}}\]

\end{enumerate}
